% !TEX program = lualatex
% Prueba breve: teoría elemental de conjuntos

\documentclass[11pt,a4paper]{article}

\usepackage{fontspec}
\usepackage[spanish,es-nodecimaldot]{babel}
\usepackage[a4paper,margin=2.2cm,headheight=15pt]{geometry}
\usepackage{amsmath,amssymb,mathtools}
\usepackage{enumitem}
\usepackage{tabularx}
\usepackage{xcolor}
\usepackage{fancyhdr}
\usepackage{lastpage}
\usepackage{hyperref}

\setmainfont{Latin Modern Roman}
\setsansfont{Latin Modern Sans}
\setmonofont{Latin Modern Mono}

\definecolor{FdMBlue}{RGB}{20,55,100}
\definecolor{FdMLightBlue}{RGB}{235,241,248}
\definecolor{FdMGray}{RGB}{90,96,104}

\hypersetup{
  colorlinks=true,
  linkcolor=FdMBlue,
  urlcolor=FdMBlue,
  pdftitle={Prueba breve: teoría elemental de conjuntos},
  pdfauthor={Fundamentos de las Matemáticas I}
}

\pagestyle{fancy}
\fancyhf{}
\lhead{Fundamentos de las Matemáticas I}
\rhead{Prueba breve: conjuntos}
\cfoot{\thepage\ de \pageref{LastPage}}
\renewcommand{\headrulewidth}{0.4pt}

\setlist[enumerate]{leftmargin=2.4em,itemsep=0.28em,topsep=0.35em}
\setlength{\parindent}{0pt}

\newcounter{pregunta}
\newenvironment{pregunta}[2][]{%
  \refstepcounter{pregunta}%
  \par\medskip\noindent
  {\color{FdMBlue}\sffamily\bfseries Pregunta \thepregunta%
  \if\relax\detokenize{#1}\relax\else\ --- #1\fi.}\hfill
  {\sffamily\bfseries[#2 puntos]}\par\smallskip
}{\par\medskip}

\newcommand{\espaciorespuesta}[1]{%
  \par\vspace{#1}\noindent
  \textcolor{FdMGray!45}{\rule{\linewidth}{0.3pt}}\par
}

\newcommand{\cajainstrucciones}[1]{%
  \begin{center}
    \colorbox{FdMLightBlue}{%
      \begin{minipage}{0.92\linewidth}
        \small #1
      \end{minipage}}
  \end{center}
}

\begin{document}

\begin{center}
  {\color{FdMBlue}\sffamily\bfseries\LARGE Prueba breve}\par
  \vspace{0.15em}
  {\sffamily\large Teoría elemental de conjuntos}\par
  \vspace{0.55em}
  \rule{\linewidth}{0.6pt}
\end{center}

\begin{tabularx}{\linewidth}{@{}lX@{\hspace{1.2em}}lX@{}}
  \textbf{Nombre y apellidos:} & \hrulefill &
  \textbf{Fecha:} & \hrulefill
\end{tabularx}

\vspace{0.6em}
\begin{tabularx}{\linewidth}{@{}lXlX@{}}
  \textbf{Grupo:} & \hrulefill &
  \textbf{Calificación:} & \hrulefill\ /\ 10
\end{tabularx}

\cajainstrucciones{
  \textbf{Duración: 30 minutos.}
  Justifica las respuestas. En las igualdades de conjuntos debe indicarse el
  razonamiento utilizado; una respuesta sin explicación puede recibir solo
  puntuación parcial.
}

\begin{pregunta}[Pertenencia e inclusión]{2}
Sean
\[
  A=\{1,2,\{1,2\},\varnothing\},
  \qquad B=\{1,2\}.
\]
Indica si cada afirmación es verdadera o falsa y justifica brevemente:
\begin{enumerate}[label=\alph*)]
  \item \(B\in A\).
  \item \(A\subseteq B\).
  \item \(\varnothing\in A\).
  \item \(\varnothing\in B\).
  \item \(\varnothing\subseteq B\).
\end{enumerate}
\espaciorespuesta{3.0cm}
\end{pregunta}

\begin{pregunta}[Operaciones con conjuntos]{2.5}
En el universo \(U=\{1,2,3,4,5,6,7,8\}\), sean
\[
  C=\{1,2,4,6\},
  \qquad D=\{2,3,6,7\}.
\]
Calcula:
\[
\begin{array}{lll}
\text{a) }C\cup D, &
\text{b) }C\cap D, &
\text{c) }C\setminus D,\\[1mm]
\text{d) }C^c, &
\text{e) }(C\cup D)^c. &
\end{array}
\]
Comprueba en este ejemplo que \((C\cup D)^c=C^c\cap D^c\).
\espaciorespuesta{3.0cm}
\end{pregunta}

\newpage

\begin{pregunta}[Producto cartesiano y conjunto de partes]{2.5}
Sean \(E=\{1,2\}\) y \(F=\{a,b\}\).
\begin{enumerate}[label=\alph*)]
  \item Escribe por extensión \(E\times F\).
  \item Calcula \(\mathcal P(E)\).
  \item Decide, justificando cada respuesta, si son verdaderas:
        \[
          \{1\}\in\mathcal P(E),\qquad
          1\in\mathcal P(E),\qquad
          E\in\mathcal P(E).
        \]
  \item Indica \(|E\times F|\) y \(|\mathcal P(E\times F)|\).
\end{enumerate}
\espaciorespuesta{6.2cm}
\end{pregunta}

\begin{pregunta}[Demostración de una identidad]{3}
Sean \(A,B,C\) subconjuntos de un universo \(U\). Demuestra, mediante
equivalencias de pertenencia o por doble inclusión, que
\[
  A\setminus(B\cup C)
  =(A\setminus B)\cap(A\setminus C).
\]
Debes señalar explícitamente dónde se utiliza la definición de diferencia.
\espaciorespuesta{8.0cm}
\end{pregunta}

\end{document}
